%! Multiplying \& Dividing Radicals; Rationalizing — Unit 6, Lesson 6.3 — Homework (Answer Key)
\documentclass[10pt]{article}
\usepackage{algebra2-article}
\usepackage{algebra2-key}   % pulls in -boxes; do NOT also load -boxes
\newcommand{\TallMath}[1]{$\displaystyle #1\rule[-1.4em]{0pt}{3.2em}$}
% In-formula answer slot (\ans is text-mode and must never appear inside $...$).
\newcommand{\mans}[1]{{\color{keyred}\mathbf{#1}}}

\begin{document}

\pageheader{Unit 6, Lesson 6.3}{Homework --- Answer Key}
\namedateperiod

\vspace{0.10in}

\begin{notesbox}{Practice}
Assume all variables represent \textbf{positive} real numbers. Leave no radical in any denominator.
No calculator except where a problem asks for a decimal.

\begin{enumerate}[leftmargin=1.9em, itemsep=11pt]
  \item \textbf{Multiply and simplify.} Remember: the product is where a perfect power often
        appears out of nowhere, so finish every one.
    \begin{center}
    \renewcommand{\arraystretch}{1.9}
    \begin{tabularx}{\linewidth}{@{}X X X X@{}}
    (a) \; $\sqrt5\cdot\sqrt{20}=$ \ans{$10$} & (b) \; $\sqrt6\cdot\sqrt{14}=$ \ans{$2\sqrt{21}$} &
    (c) \; $4\sqrt3\cdot2\sqrt6=$ \ans{$24\sqrt2$} & (d) \; $\left(3\sqrt7\right)^{2}=$ \ans{$63$} \\
    (e) \; $\sqrt[3]{4}\cdot\sqrt[3]{10}=$ \ans{$2\sqrt[3]{5}$} &
    (f) \; $\sqrt{6x}\cdot\sqrt{3x^{3}}=$ \ans{$3x^{2}\sqrt2$} &
    (g) \; $\sqrt{10}\cdot\sqrt{10}=$ \ans{$10$} &
    (h) \; $2\sqrt{12}\cdot\sqrt{27}=$ \ans{$36$} \\
    \end{tabularx}
    \end{center}
    \emph{Work:} (a) $\sqrt{100}$. \; (b) $\sqrt{84}=\sqrt{4\cdot21}$. \;
    (c) $8\sqrt{18}=8\cdot3\sqrt2$. \; (d) $9\cdot7$ --- squaring erases the radical. \;
    (e) $\sqrt[3]{40}=\sqrt[3]{8\cdot5}$. \; (f) $\sqrt{18x^{4}}=\sqrt{9x^{4}}\cdot\sqrt2$. \;
    (g) $\sqrt{100}$ again --- a square root times itself is the radicand. \;
    (h) $2\sqrt{324}=2\cdot18$.

  \item \textbf{Expand and simplify.} Distribute or FOIL, then collect like radicals.
    \begin{center}
    \renewcommand{\arraystretch}{2.1}
    \begin{tabularx}{\linewidth}{@{}X X@{}}
    (a) \; $\sqrt2\left(\sqrt8+\sqrt{18}\right)=$ \ans{$10$} &
    (b) \; $\sqrt5\left(3-\sqrt{15}\right)=$ \ans{$3\sqrt5-5\sqrt3$} \\
    (c) \; $\left(1+\sqrt3\right)\left(4-\sqrt3\right)=$ \ans{$1+3\sqrt3$} &
    (d) \; $\left(\sqrt6-\sqrt2\right)^{2}=$ \ans{$8-4\sqrt3$} \\
    (e) \; $\left(5+\sqrt7\right)\left(5-\sqrt7\right)=$ \ans{$18$} &
    (f) \; $\left(2\sqrt3+\sqrt5\right)\left(2\sqrt3-\sqrt5\right)=$ \ans{$7$} \\
    \end{tabularx}
    \end{center}
    \emph{Work:} (a) $\sqrt{16}+\sqrt{36}=4+6$. \; (b) $3\sqrt5-\sqrt{75}=3\sqrt5-5\sqrt3$. \;
    (c) $4-\sqrt3+4\sqrt3-3$. \; (d) $6-2\sqrt{12}+2=8-4\sqrt3$. \; (e) $25-7$. \;
    (f) $\left(2\sqrt3\right)^{2}-\left(\sqrt5\right)^{2}=12-5$.

    \vspace{2pt}
    Two of the six are \textbf{conjugate pairs}, and you could tell before multiplying. Which two?
    \ans{(e) and (f)} \; A \emph{third} problem also came out with no radical, for a completely
    different reason. Which one, and what was the reason? \ansline{\textbf{(a).} Nothing cancelled
    there --- it is not a conjugate pair at all. Distributing produced $\sqrt{16}$ and $\sqrt{36}$,
    and both radicands happened to be \emph{perfect squares}, so each term came out whole on its
    own. (e) and (f) are rational because the middle terms cancel; (a) is rational by luck of the
    numbers.}

  \item \textbf{Divide, and rationalize where you must.}
    \begin{center}
    \renewcommand{\arraystretch}{2.3}
    \begin{tabularx}{\linewidth}{@{}X X@{}}
    (a) \; $\dfrac{\sqrt{72}}{\sqrt{2}}=$ \ans{$6$} &
    (b) \; $\dfrac{\sqrt{45x^{3}}}{\sqrt{5x}}=$ \ans{$3x$} \\
    (c) \; $\dfrac{7}{\sqrt{2}}=$ \ans{$\dfrac{7\sqrt2}{2}$} &
    (d) \; $\dfrac{\sqrt{3}}{\sqrt{8}}=$ \ans{$\dfrac{\sqrt6}{4}$} \\
    (e) \; $\sqrt{\dfrac{2}{7}}=$ \ans{$\dfrac{\sqrt{14}}{7}$} &
    (f) \; $\dfrac{10}{4+\sqrt{6}}=$ \ans{$4-\sqrt6$} \\
    (g) \; $\dfrac{\sqrt{2}}{\sqrt{5}+\sqrt{3}}=$ \ans{$\dfrac{\sqrt{10}-\sqrt6}{2}$} &
    (h) \; $\dfrac{6}{\sqrt{3}-3}=$ \ans{$-\sqrt3-3$} \\
    \end{tabularx}
    \end{center}
    \emph{Work:} (a) $\sqrt{36}$. \; (b) $\sqrt{9x^{2}}$. \; (d) simplify $\sqrt8=2\sqrt2$ first,
    then $\frac{\sqrt3}{2\sqrt2}\cdot\frac{\sqrt2}{\sqrt2}=\frac{\sqrt6}{4}$. \;
    (f) conjugate $4-\sqrt6$; denominator $16-6=10$, which cancels the $10$ on top entirely. \;
    (g) conjugate $\sqrt5-\sqrt3$; denominator $5-3=2$. \;
    (h) conjugate $\sqrt3+3$; denominator $3-9=-6$, so
    $\frac{6\left(\sqrt3+3\right)}{-6}=-\left(\sqrt3+3\right)$. \emph{The sign in (h) is the whole
    point} --- a conjugate product can be negative, and the minus sign belongs to the entire
    numerator.

  \item \textbf{Two questions about why this works.}
    \begin{enumerate}[label=(\alph*), leftmargin=1.8em, itemsep=6pt]
      \item A classmate claims $\left(\sqrt3+\sqrt5\right)^{2}=3+5=8$. Use decimals to show this is
            wrong, then expand it correctly and name the term they threw away.
            \ansline{$\sqrt3+\sqrt5\approx1.732+2.236=3.968$, and $3.968^{2}\approx15.75$ --- not
            $8$. Correctly: $\left(\sqrt3+\sqrt5\right)^{2}=3+\sqrt{15}+\sqrt{15}+5=8+2\sqrt{15}$,
            and $8+2(3.873)\approx15.75$, which matches.\\[3pt]
            The discarded term is the \textbf{middle term} $2\sqrt{15}$ --- the two ``outer and
            inner'' products of FOIL. It is the same error as $(x+4)^{2}=x^{2}+16$: squaring a sum
            is never just squaring the pieces.}
      \item To rationalize $\dfrac{1}{2+\sqrt3}$, why does multiplying top and bottom by
            $2+\sqrt3$ fail while $2-\sqrt3$ succeeds? And why is either one \emph{allowed} in the
            first place --- what keeps the value of the expression from changing?
            \ansline{Multiplying by $2+\sqrt3$ gives $\left(2+\sqrt3\right)^{2}=7+4\sqrt3$ --- both
            middle products are $+2\sqrt3$, so they add and the radical survives. Multiplying by
            the \emph{conjugate} $2-\sqrt3$ gives $4-3=1$: the middle products $-2\sqrt3$ and
            $+2\sqrt3$ cancel, and each remaining term is a square, so no radical is left. (The
            answer is $2-\sqrt3$.)\\[3pt]
            Either is \emph{allowed} because in both cases the numerator and denominator are
            multiplied by the \emph{same} expression --- that is multiplying by a fraction equal to
            $1$, which changes an expression's form but never its value. Only one of the two is
            \emph{useful}.}
    \end{enumerate}

  \item \textbf{Two roads, one answer (connecting to Lesson 6.1).} Rationalize
        $\dfrac{1}{\sqrt[3]{5}}$ two ways.
    \begin{enumerate}[label=(\alph*), leftmargin=1.8em, itemsep=5pt]
      \item \emph{With radicals.} You need \ans{$3$} copies of $5$ under the cube root, so
            multiply top and bottom by \ans{$\sqrt[3]{25}$}: \;
            $\dfrac{1}{\sqrt[3]5}=$ \ans{$\dfrac{\sqrt[3]{25}}{5}$}
      \item \emph{With exponents.} $\dfrac{1}{\sqrt[3]5}=5^{\mans{-1/3}}
            =\dfrac{5^{\mans{2/3}}}{5}=$ \ans{$\dfrac{\sqrt[3]{25}}{5}$}
      \item Why must the two roads agree? \ansline{They are the same operation written twice.
            Multiplying by $\sqrt[3]{25}=5^{2/3}$ raises the denominator's exponent from
            $\frac13$ to $\frac13+\frac23=1$, which is exactly what ``complete the group of three''
            means. In exponent language, rationalizing is just turning the negative exponent
            $5^{-1/3}$ into a positive one by borrowing a whole $5$: $5^{-1/3}=5^{2/3-1}
            =\frac{5^{2/3}}{5}$.}
    \end{enumerate}
\end{enumerate}
\end{notesbox}

\begin{extensionbox}
\textbf{The golden ratio: the number that is its own reciprocal, plus one.} The
\textbf{golden ratio} is
\[
  \varphi = \frac{1+\sqrt{5}}{2}.
\]
It shows up wherever a shape can be divided so that the part relates to the whole the way the
smaller piece relates to the part. Everything interesting about $\varphi$ is invisible until you
rationalize.

\begin{center}
\begin{tikzpicture}[scale=2.4, line join=round]
  % golden rectangle, width phi = 1.618, height 1
  \draw[thick, forest, fill=forestbg] (0,0) rectangle (1.618,1);
  \draw[thick, goldacc] (1,0) -- (1,1);
  \node[charcoal] at (0.5,0.55) {\footnotesize square};
  \node[charcoal] at (0.5,0.38) {\footnotesize $1\times 1$};
  \node[charcoal, align=center] at (1.309,0.5) {\footnotesize what's\\[-2pt] \footnotesize left};
  % dimension labels
  \draw[<->, slate] (0,-0.13) -- (1.618,-0.13);
  \node[below, charcoal] at (0.809,-0.13) {\footnotesize $\varphi$};
  \draw[<->, slate] (-0.13,0) -- (-0.13,1);
  \node[left, charcoal] at (-0.13,0.5) {\footnotesize $1$};
  \draw[<->, slate] (1,1.13) -- (1.618,1.13);
  \node[above, charcoal] at (1.309,1.13) {\footnotesize $\varphi-1$};
\end{tikzpicture}
\end{center}

\begin{enumerate}[label=(\alph*), leftmargin=1.8em, itemsep=5pt]
  \item Using $\sqrt5\approx2.236068$, fill in the table to four decimal places.
        \begin{center}
        \renewcommand{\arraystretch}{1.4}
        \begin{tabular}{|c|c|c|}
        \hline
        $\varphi$ & $\dfrac{1}{\varphi}$ & $\varphi^{2}$ \\ \hline
        \ans{$1.6180$} & \ans{$0.6180$} & \ans{$2.6180$} \\ \hline
        \end{tabular}
        \end{center}
        Something very strange is going on with those three decimals. Describe it in one sentence.
        \ansline{All three have the \emph{identical} decimal part $.6180$: the reciprocal is exactly
        one \emph{less} than $\varphi$, and the square is exactly one \emph{more}. No other positive
        number does this.}
  \item \textbf{Prove the first strange thing:} $\dfrac{1}{\varphi}=\varphi-1$. Start from
        $\dfrac{1}{\varphi}=\dfrac{2}{1+\sqrt5}$, rationalize with the conjugate, and compare your
        result to $\varphi-1$. \ansline{$\dfrac{2}{1+\sqrt5}\cdot\dfrac{1-\sqrt5}{1-\sqrt5}
        =\dfrac{2\left(1-\sqrt5\right)}{1-5}=\dfrac{2\left(1-\sqrt5\right)}{-4}
        =\dfrac{\sqrt5-1}{2}$.\\[3pt]
        And $\varphi-1=\dfrac{1+\sqrt5}{2}-\dfrac{2}{2}=\dfrac{\sqrt5-1}{2}$. Identical.}
  \item \textbf{Prove the second:} $\varphi^{2}=\varphi+1$. Square $\dfrac{1+\sqrt5}{2}$, simplify,
        and compare to $\varphi+1$ written over a common denominator.
        \ansline{$\varphi^{2}=\dfrac{\left(1+\sqrt5\right)^{2}}{4}=\dfrac{1+2\sqrt5+5}{4}
        =\dfrac{6+2\sqrt5}{4}=\dfrac{3+\sqrt5}{2}$ --- note the middle term $2\sqrt5$, which is
        exactly what item 4(a) warned you not to drop.\\[3pt]
        And $\varphi+1=\dfrac{1+\sqrt5}{2}+\dfrac{2}{2}=\dfrac{3+\sqrt5}{2}$. Identical.}
  \item \textbf{The rectangle.} The figure above is $\varphi$ wide and $1$ tall. Cut off the
        $1\times1$ square. The leftover rectangle is $1$ tall and $\varphi-1$ wide, so its
        \emph{long side divided by its short side} is $\dfrac{1}{\varphi-1}$. Rationalize that
        expression and show it equals $\varphi$ --- the leftover rectangle has exactly the same
        proportions as the original. \ansline{From (b), $\varphi-1=\dfrac{\sqrt5-1}{2}$, so
        $\dfrac{1}{\varphi-1}=\dfrac{2}{\sqrt5-1}$. Multiply by the conjugate:
        $\dfrac{2}{\sqrt5-1}\cdot\dfrac{\sqrt5+1}{\sqrt5+1}
        =\dfrac{2\left(\sqrt5+1\right)}{5-1}=\dfrac{\sqrt5+1}{2}=\varphi$.\\[3pt]
        So the leftover rectangle's long-to-short ratio is $\varphi$ as well --- it is another
        golden rectangle, just smaller and turned on its side.}
  \item What does part (d) mean about repeating the cut on the leftover rectangle, and then again
        on \emph{its} leftover? Why is this one of the few facts about $\varphi$ you could not have
        seen without rationalizing? \ansline{The cut can be repeated \textbf{forever}. Every
        leftover is golden, so every leftover can have a square removed to produce yet another
        golden rectangle --- an infinite nested spiral of them, which is why this shape appears in
        art and design.\\[3pt]
        You could not see it in the original form because $\dfrac{1}{\varphi-1}$ has a radical
        buried in its denominator, and in that shape it looks like nothing in particular.
        Rationalizing turns it into $\dfrac{1+\sqrt5}{2}$, and only then is it recognizably
        $\varphi$ again. Rationalizing did not compute anything new --- it made an equality
        \emph{visible}.}
\end{enumerate}
\end{extensionbox}

\begin{spiralbox}
\textbf{Coming next --- Lesson 6.4: Graphing Radical Functions \& Transformations.} The algebra of
radicals is now finished: you can simplify them, add them, multiply them, and clear them out of a
denominator. Next you go back to the \emph{pictures} from Lesson 6.0 and learn to move them around:
$y=a\sqrt{x-h}+k$ and $y=a\sqrt[3]{x-h}+k$, where the \textbf{endpoint} is the anchor that every
transformation moves, and $a$, $h$, and $k$ do exactly what they did to parabolas back in Unit 3.
You will read a graph and write its equation, and write an equation and predict its graph.
\end{spiralbox}

\begin{teachernote}
Items 1--3 are the routine load, about twenty-five minutes. Two spot-checks are worth your time.
\textbf{Item 2(a)} is a trap for anyone who distributes without finishing: $\sqrt{16}+\sqrt{36}$
looks like an answer and is not. \textbf{Item 3(h)} is the one problem on the page with a
\emph{negative} conjugate product ($3-9=-6$); students who write $\sqrt3+3$ instead of
$-\sqrt3-3$ have done everything right except distribute the minus sign over both terms, which is
the Unit 4 leading-minus error resurfacing.

\textbf{Item 2's follow-up question is the best diagnostic on the page.} A student who lists (a)
among the conjugate pairs has learned ``no radical in the answer $\Rightarrow$ conjugates'' as a
pattern rather than a mechanism. The distinction matters: (e) and (f) are rational \emph{by
structure} and can be predicted before any multiplying; (a) is rational because two particular
radicands happened to be perfect squares.

\textbf{Item 4(b) is the conceptual spine} and should be graded on the explanation, not the answer.
Two separate things have to be said --- \emph{why the conjugate works} (middle terms cancel) and
\emph{why multiplying by it is legal at all} (it is a form of $1$). Students who can say only the
first will treat rationalizing as a ritual; they are the ones who will not think to use it in
Lesson 6.5 when it is not asked for by name.

\textbf{Item 5} closes the same kind of loop Lesson 6.2's homework closed: a procedure and its
rational-exponent explanation, meeting in the middle. Accept any wording that connects
``complete the group of three'' to $\frac13+\frac23=1$.

The \textbf{Extension} is the golden ratio, and its shape deliberately matches the unit's earlier
models --- a table of exact values, a surprising pattern, a structural rewrite. But its point is
different from the pendulum's or Kepler's: here rationalizing is not a tidying-up step, it is the
\emph{only} way to see that $\frac{1}{\varphi-1}$ and $\varphi$ are the same number. Part (a) is
worth the whole assignment on its own --- students genuinely do not expect $1.6180$, $0.6180$, and
$2.6180$ --- and parts (b)--(d) turn that observation into proof. Part (c) quietly re-tests the
middle-term error from item 4(a): a student who squares $1+\sqrt5$ as $1+5$ gets
$\varphi^{2}=\frac{6}{4}$ and the pattern collapses, which is a self-correcting check.
\end{teachernote}

\end{document}
